\section{Conclusion}
In this paper, we presented \textbf{\our}, a novel framework that shifts the paradigm of generative world modeling from 2D pixel sequences to consistent 4D scene evolution. By introducing a lightweight yet expressive representation consisting of synchronized RGB, depth, and optical flow (RGB-DF), we effectively bridge the gap between the scalability of video diffusion models and the geometric rigor required for robotic manipulation. Our core contributions include the development of a tri-branch architecture that ensures cross-modal consistency through mutual feature interactions, and the curation of \textbf{Rynn4DDataset 1.0}, the large-scale hybrid dataset specifically designed for training 4D generative models. Furthermore, we demonstrated that \textbf{\our-Policy} can effectively leverage these predictive 4D representations as an implicit world model, enabling high-frequency, closed-loop robotic control that outperforms existing 2D-based baselines. Extensive experiments show that our approach not only generates high-fidelity 4D futures but also significantly enhances the success rate and precision of downstream manipulation tasks. We believe \our provides a promising foundation for building general-purpose embodied intelligence capable of understanding and interacting with the complex 3D world.

\noindent
\textbf{Limitation.} 
Despite the reactive capabilities of \our-Policy, our framework has several limitations. First, the 4D sequence generation relies on a diffusion denoising process, which introduces computational overhead. Currently, our implementation achieves an effective control frequency of approximately 9 Hz on an NVIDIA RTX 5090 GPU; while sufficient for many tasks, this latency remains a bottleneck for ultra-high-frequency control. Second, our model is primarily optimized for egocentric perspectives. Extending 4D spatio-temporal consistency to multi-view systems or collaborative multi-robot setups remains an open challenge for future research.
